\documentclass[aspectratio=169,10pt]{beamer}
% ═══════════════════════════════════════════════════════════════════════════════
% SHARED PREAMBLE — Foundation Models Course (HIT)
% To be \input by each lecture file AFTER \documentclass
% ═══════════════════════════════════════════════════════════════════════════════

% ─────────────────────────────────────────────────────────────────────────────
% BEAMER THEME & BASIC SETUP
% ─────────────────────────────────────────────────────────────────────────────
\usetheme{Madrid}
\usecolortheme{default}

% Remove navigation symbols
\beamertemplatenavigationsymbolsempty

% ─────────────────────────────────────────────────────────────────────────────
% COLORS — HIT Branding
% ─────────────────────────────────────────────────────────────────────────────
\definecolor{hitblue}{RGB}{0,51,102}        % Primary: HIT Navy
\definecolor{hitgold}{RGB}{192,148,0}       % Accent: HIT Gold
\definecolor{hitlight}{RGB}{230,240,250}    % Light background

% Override Beamer palette
\setbeamercolor{primary}{fg=white,bg=hitblue}
\setbeamercolor{secondary}{fg=hitblue,bg=hitlight}
\setbeamercolor{structure}{fg=hitblue}
\setbeamercolor{alerted text}{fg=hitgold}

% ─────────────────────────────────────────────────────────────────────────────
% PACKAGES
% ─────────────────────────────────────────────────────────────────────────────
\usepackage{amsmath}
\usepackage{amssymb}
\usepackage{mathtools}
\usepackage{booktabs}
\usepackage{graphicx}
\usepackage{hyperref}
\usepackage{tikz}
\usepackage{pgfplots}
\usepackage{tcolorbox}
\usepackage{listings}
\usepackage{xcolor}
\usepackage{algorithm}
\usepackage{algpseudocode}

% ─────────────────────────────────────────────────────────────────────────────
% TikZ LIBRARIES
% ─────────────────────────────────────────────────────────────────────────────
\usetikzlibrary{shapes,arrows,positioning,calc,chains,scopes}
\pgfplotsset{compat=1.17}

% ─────────────────────────────────────────────────────────────────────────────
% CODE LISTINGS
% ─────────────────────────────────────────────────────────────────────────────
\lstset{
  basicstyle=\ttfamily\small,
  breaklines=true,
  commentstyle=\color{gray},
  keywordstyle=\color{hitblue},
  stringstyle=\color{hitgold},
  showstringspaces=false,
  backgroundcolor=\color{hitlight},
  frame=single,
  rulecolor=\color{hitblue}
}

% ─────────────────────────────────────────────────────────────────────────────
% TCOLORBOX STYLES
% ─────────────────────────────────────────────────────────────────────────────
\tcbset{
  hitbox/.style={
    colback=hitlight,
    colframe=hitblue,
    fonttitle=\bfseries,
    boxrule=1pt
  },
  alertbox/.style={
    colback=yellow!10,
    colframe=hitgold,
    fonttitle=\bfseries,
    boxrule=1.5pt
  }
}

% ─────────────────────────────────────────────────────────────────────────────
% CUSTOM COMMANDS
% ─────────────────────────────────────────────────────────────────────────────

% Placeholder for figures
\newcommand{\placeholder}[3]{
  \begin{center}
    \fbox{\begin{minipage}[c][#2][c]{#1}
      \centering\small\color{gray}[Figure: #3]
    \end{minipage}}
  \end{center}
}

% Section divider frame
\newcommand{\sectionframe}[2]{
  \begin{frame}[plain]
    \begin{tikzpicture}[remember picture,overlay]
      \fill[hitblue] (current page.south west) rectangle (current page.north east);
      \node[anchor=center, white, font=\Large\bfseries] at (current page.center) {
        \begin{tabular}{c}
          #1 \\[0.3cm]
          {\small\color{hitgold}#2}
        \end{tabular}
      };
    \end{tikzpicture}
  \end{frame}
}

% ─────────────────────────────────────────────────────────────────────────────
% LOAD CUSTOM MACROS
% ─────────────────────────────────────────────────────────────────────────────
% ═══════════════════════════════════════════════════════════════════════════════
% SHARED MACROS — Mathematical Notation for Foundation Models Course
% ═══════════════════════════════════════════════════════════════════════════════

% ─────────────────────────────────────────────────────────────────────────────
% ACTIVATION & LAYER FUNCTIONS
% ─────────────────────────────────────────────────────────────────────────────
\newcommand{\softmax}{\mathrm{softmax}}
\newcommand{\sigmoid}{\mathrm{sigmoid}}
\newcommand{\relu}{\mathrm{ReLU}}
\newcommand{\gelu}{\mathrm{GELU}}
\newcommand{\LayerNorm}{\mathrm{LayerNorm}}
\newcommand{\FFN}{\mathrm{FFN}}
\newcommand{\MHA}{\mathrm{MHA}}
\newcommand{\Attn}{\mathrm{Attention}}

% ─────────────────────────────────────────────────────────────────────────────
% ATTENTION COMPONENTS
% ─────────────────────────────────────────────────────────────────────────────
\newcommand{\query}{Q}
\newcommand{\key}{K}
\newcommand{\val}{V}
\newcommand{\dmodel}{d_{\text{model}}}
\newcommand{\dk}{d_k}
\newcommand{\nheads}{h}

% Scaled dot-product attention formula
\newcommand{\SDPA}{
  \Attn(\query, \key, \val) = \softmax\left(\frac{\query \key^T}{\sqrt{\dk}}\right) \val
}

\newcommand{\CrossAttn}{\text{CrossAttention}}

% ─────────────────────────────────────────────────────────────────────────────
% PROBABILITY & EXPECTATION
% ─────────────────────────────────────────────────────────────────────────────
\newcommand{\E}{\mathbb{E}}
\newcommand{\Prob}{\mathbb{P}}
\newcommand{\KL}{\mathrm{KL}}
\newcommand{\ELBO}{\mathrm{ELBO}}
\newcommand{\given}{\mid}

% ─────────────────────────────────────────────────────────────────────────────
% NORMS & OPERATIONS
% ─────────────────────────────────────────────────────────────────────────────
\newcommand{\norm}[1]{\left\lVert #1 \right\rVert}
\newcommand{\abs}[1]{\left| #1 \right|}
\newcommand{\inner}[2]{\langle #1, #2 \rangle}
\newcommand{\T}{^{\top}}

% ─────────────────────────────────────────────────────────────────────────────
% VECTORS & MATRICES
% ─────────────────────────────────────────────────────────────────────────────
\newcommand{\bvec}[1]{\mathbf{#1}}
\newcommand{\mat}[1]{\mathbf{#1}}
\newcommand{\iid}{\stackrel{\text{i.i.d.}}{\sim}}
\newcommand{\defeq}{\coloneq}

% ─────────────────────────────────────────────────────────────────────────────
% LANGUAGE MODELING
% ─────────────────────────────────────────────────────────────────────────────
\newcommand{\vocab}{\mathcal{V}}
\newcommand{\context}{\text{context}}
\newcommand{\token}{t}
\newcommand{\seq}{x}
\newcommand{\ptheta}{p_\theta}
\newcommand{\pref}{p_{\text{ref}}}

% ─────────────────────────────────────────────────────────────────────────────
% RLHF & DPO
% ─────────────────────────────────────────────────────────────────────────────
\newcommand{\piref}{\pi_{\text{ref}}}
\newcommand{\pitheta}{\pi_\theta}
\newcommand{\yw}{y_w}
\newcommand{\yl}{y_l}
\newcommand{\DPOloss}{\mathcal{L}_{\mathrm{DPO}}}

% ─────────────────────────────────────────────────────────────────────────────
% RETRIEVAL & RAG
% ─────────────────────────────────────────────────────────────────────────────
\newcommand{\docs}{\mathcal{D}}
\newcommand{\qvec}{q}
\newcommand{\dvec}{d}

% ─────────────────────────────────────────────────────────────────────────────
% AGENTS & REASONING
% ─────────────────────────────────────────────────────────────────────────────
\newcommand{\obs}{o}
\newcommand{\act}{a}
\newcommand{\thought}{t}
\newcommand{\tools}{\mathcal{T}}
\newcommand{\history}{h}
\newcommand{\concat}{\oplus}

% ─────────────────────────────────────────────────────────────────────────────
% SETS & LOGIC
% ─────────────────────────────────────────────────────────────────────────────
\newcommand{\R}{\mathbb{R}}
\newcommand{\N}{\mathbb{N}}
\newcommand{\Z}{\mathbb{Z}}

\endinput


% ─────────────────────────────────────────────────────────────────────────────
% TITLE & AUTHOR (can be overridden in individual lectures)
% ─────────────────────────────────────────────────────────────────────────────
\author[D. Lederman]{Dror Lederman}
\institute{Holon Institute of Technology (HIT) \\ Data Science Department}
\date{\today}

% ─────────────────────────────────────────────────────────────────────────────
% FOOTER & HEADER
% ─────────────────────────────────────────────────────────────────────────────
\setbeamertemplate{footline}{
  \leavevmode%
  \hbox{%
    \begin{beamercolorbox}[wd=0.33\paperwidth,ht=2.25ex,dp=1ex,center]{author in head/foot}%
      \usebeamerfont{author in head/foot}\insertshortauthor
    \end{beamercolorbox}%
    \begin{beamercolorbox}[wd=0.34\paperwidth,ht=2.25ex,dp=1ex,center]{title in head/foot}%
      \usebeamerfont{title in head/foot}\insertshorttitle
    \end{beamercolorbox}%
    \begin{beamercolorbox}[wd=0.33\paperwidth,ht=2.25ex,dp=1ex,right]{frame number}%
      \usebeamerfont{frame number}\insertframenumber{} / \inserttotalframenumber\hspace*{2ex}
    \end{beamercolorbox}%
  }%
  \vskip0pt%
}

\endinput


\title{Week 6: Retrieval-Augmented Generation}
\subtitle{Foundation Models and Agentic AI}
\author{Lecturer Name}
\institute{Holon Institute of Technology (HIT)}
\date{Semester, 2025}

\begin{document}

\begin{frame}
  \titlepage
\end{frame}

\begin{frame}{Learning Objectives}
  After this lecture you will be able to:
  \begin{itemize}
    \item Describe the \textbf{RAG architecture} and why it reduces hallucination.
    \item Apply \textbf{chunking strategies} and understand their impact on retrieval quality.
    \item Implement \textbf{query rewriting} and \textbf{reranking} for advanced RAG.
    \item Explain \textbf{Self-RAG}: selective retrieval with self-critique.
    \item Understand \textbf{Graph RAG} for multi-hop and global summarization queries.
    \item Evaluate RAG systems using \textbf{RAGAS} and related frameworks.
  \end{itemize}
\end{frame}

\section{Why RAG?}
\sectionframe{Why Retrieval-Augmented Generation?}{Grounding LLMs in external knowledge}

\begin{frame}{Limitations of Parametric LLMs}
  \begin{tcolorbox}[alertbox, title=Core Problems]
    \begin{itemize}
      \item \textbf{Knowledge cutoff:} LLMs cannot answer about events after training.
      \item \textbf{Hallucination:} LLMs confabulate plausible-sounding but false facts.
      \item \textbf{Opacity:} Cannot cite sources; hard to audit answers.
      \item \textbf{Scale cost:} Updating knowledge requires expensive retraining.
    \end{itemize}
  \end{tcolorbox}
  \vspace{0.4em}
  \textbf{RAG solution:} At inference time, retrieve relevant documents and provide them as
  context alongside the query. The LLM reads documents; it doesn't need to memorize facts.
  \note{RAG was independently motivated by QA systems (open-book exam analogy) and production search systems.}
\end{frame}

\begin{frame}{RAG Architecture Overview}
  % tikz/rag_pipeline.tex
% RAG pipeline: offline indexing path + online query path
% Usage: % tikz/rag_pipeline.tex
% RAG pipeline: offline indexing path + online query path
% Usage: % tikz/rag_pipeline.tex
% RAG pipeline: offline indexing path + online query path
% Usage: \input{../../tikz/rag_pipeline}
\begin{tikzpicture}[
  node distance=0.6cm and 0.8cm,
  box/.style={rectangle, draw=hitblue, fill=hitlightblue, rounded corners=3pt,
              minimum width=1.9cm, minimum height=0.6cm, font=\scriptsize, align=center},
  db/.style={cylinder, draw=hitgold, fill=hitgold!15, shape border rotate=90,
             minimum width=1.5cm, minimum height=0.8cm, font=\scriptsize, align=center},
  io/.style={rectangle, draw=hitgold, fill=hitgold!15, rounded corners=3pt,
             minimum width=1.9cm, minimum height=0.6cm, font=\scriptsize, align=center},
  arr/.style={->, >=Stealth, thick, hitblue},
  darr/.style={->, >=Stealth, thick, hitgold, dashed},
  lbl/.style={font=\tiny, gray},
]

%% ── Offline Indexing Path (top row) ──────────────────────────────────────────
\node[io] (docs) {Documents};
\node[box, right=0.7cm of docs] (chunk) {Chunker};
\node[box, right=0.7cm of chunk] (enc_off) {Encoder\\(Embed)};
\node[db, right=0.7cm of enc_off] (vdb) {Vector\\DB};

\draw[darr] (docs) -- (chunk);
\draw[darr] (chunk) -- (enc_off);
\draw[darr] (enc_off) -- node[above, lbl] {vectors} (vdb);

% Offline label
\node[lbl, above=0.15cm of docs, font=\tiny\bfseries, gray] {Offline Indexing};

%% ── Online Query Path (bottom row) ───────────────────────────────────────────
\node[io, below=1.4cm of docs] (query) {User Query};
\node[box, right=0.7cm of query] (enc_q) {Query\\Encoder};
\node[box, right=0.7cm of enc_q] (ret) {Retriever\\(ANN)};
\node[box, right=0.7cm of ret] (rerank) {Reranker\\(opt.)};
\node[box, right=0.7cm of rerank] (gen) {LLM\\Generator};
\node[io, right=0.7cm of gen] (ans) {Answer};

\draw[arr] (query) -- (enc_q);
\draw[arr] (enc_q) -- (ret);
\draw[arr] (ret) -- (rerank);
\draw[arr] (rerank) -- (gen);
\draw[arr] (gen) -- (ans);

% Retriever pulls from VDB
\draw[darr, bend right=25] (vdb.south) to
  node[right, lbl, pos=0.5] {top-$k$} (ret.north);

% Context arrow into generator
\draw[darr, bend left=20] (rerank.south) to
  node[below, lbl, pos=0.5] {retrieved chunks} (gen.south west);

% Online label
\node[lbl, above=0.15cm of query, font=\tiny\bfseries, gray] {Online Query};

\end{tikzpicture}

\begin{tikzpicture}[
  node distance=0.6cm and 0.8cm,
  box/.style={rectangle, draw=hitblue, fill=hitlightblue, rounded corners=3pt,
              minimum width=1.9cm, minimum height=0.6cm, font=\scriptsize, align=center},
  db/.style={cylinder, draw=hitgold, fill=hitgold!15, shape border rotate=90,
             minimum width=1.5cm, minimum height=0.8cm, font=\scriptsize, align=center},
  io/.style={rectangle, draw=hitgold, fill=hitgold!15, rounded corners=3pt,
             minimum width=1.9cm, minimum height=0.6cm, font=\scriptsize, align=center},
  arr/.style={->, >=Stealth, thick, hitblue},
  darr/.style={->, >=Stealth, thick, hitgold, dashed},
  lbl/.style={font=\tiny, gray},
]

%% ── Offline Indexing Path (top row) ──────────────────────────────────────────
\node[io] (docs) {Documents};
\node[box, right=0.7cm of docs] (chunk) {Chunker};
\node[box, right=0.7cm of chunk] (enc_off) {Encoder\\(Embed)};
\node[db, right=0.7cm of enc_off] (vdb) {Vector\\DB};

\draw[darr] (docs) -- (chunk);
\draw[darr] (chunk) -- (enc_off);
\draw[darr] (enc_off) -- node[above, lbl] {vectors} (vdb);

% Offline label
\node[lbl, above=0.15cm of docs, font=\tiny\bfseries, gray] {Offline Indexing};

%% ── Online Query Path (bottom row) ───────────────────────────────────────────
\node[io, below=1.4cm of docs] (query) {User Query};
\node[box, right=0.7cm of query] (enc_q) {Query\\Encoder};
\node[box, right=0.7cm of enc_q] (ret) {Retriever\\(ANN)};
\node[box, right=0.7cm of ret] (rerank) {Reranker\\(opt.)};
\node[box, right=0.7cm of rerank] (gen) {LLM\\Generator};
\node[io, right=0.7cm of gen] (ans) {Answer};

\draw[arr] (query) -- (enc_q);
\draw[arr] (enc_q) -- (ret);
\draw[arr] (ret) -- (rerank);
\draw[arr] (rerank) -- (gen);
\draw[arr] (gen) -- (ans);

% Retriever pulls from VDB
\draw[darr, bend right=25] (vdb.south) to
  node[right, lbl, pos=0.5] {top-$k$} (ret.north);

% Context arrow into generator
\draw[darr, bend left=20] (rerank.south) to
  node[below, lbl, pos=0.5] {retrieved chunks} (gen.south west);

% Online label
\node[lbl, above=0.15cm of query, font=\tiny\bfseries, gray] {Online Query};

\end{tikzpicture}

\begin{tikzpicture}[
  node distance=0.6cm and 0.8cm,
  box/.style={rectangle, draw=hitblue, fill=hitlightblue, rounded corners=3pt,
              minimum width=1.9cm, minimum height=0.6cm, font=\scriptsize, align=center},
  db/.style={cylinder, draw=hitgold, fill=hitgold!15, shape border rotate=90,
             minimum width=1.5cm, minimum height=0.8cm, font=\scriptsize, align=center},
  io/.style={rectangle, draw=hitgold, fill=hitgold!15, rounded corners=3pt,
             minimum width=1.9cm, minimum height=0.6cm, font=\scriptsize, align=center},
  arr/.style={->, >=Stealth, thick, hitblue},
  darr/.style={->, >=Stealth, thick, hitgold, dashed},
  lbl/.style={font=\tiny, gray},
]

%% ── Offline Indexing Path (top row) ──────────────────────────────────────────
\node[io] (docs) {Documents};
\node[box, right=0.7cm of docs] (chunk) {Chunker};
\node[box, right=0.7cm of chunk] (enc_off) {Encoder\\(Embed)};
\node[db, right=0.7cm of enc_off] (vdb) {Vector\\DB};

\draw[darr] (docs) -- (chunk);
\draw[darr] (chunk) -- (enc_off);
\draw[darr] (enc_off) -- node[above, lbl] {vectors} (vdb);

% Offline label
\node[lbl, above=0.15cm of docs, font=\tiny\bfseries, gray] {Offline Indexing};

%% ── Online Query Path (bottom row) ───────────────────────────────────────────
\node[io, below=1.4cm of docs] (query) {User Query};
\node[box, right=0.7cm of query] (enc_q) {Query\\Encoder};
\node[box, right=0.7cm of enc_q] (ret) {Retriever\\(ANN)};
\node[box, right=0.7cm of ret] (rerank) {Reranker\\(opt.)};
\node[box, right=0.7cm of rerank] (gen) {LLM\\Generator};
\node[io, right=0.7cm of gen] (ans) {Answer};

\draw[arr] (query) -- (enc_q);
\draw[arr] (enc_q) -- (ret);
\draw[arr] (ret) -- (rerank);
\draw[arr] (rerank) -- (gen);
\draw[arr] (gen) -- (ans);

% Retriever pulls from VDB
\draw[darr, bend right=25] (vdb.south) to
  node[right, lbl, pos=0.5] {top-$k$} (ret.north);

% Context arrow into generator
\draw[darr, bend left=20] (rerank.south) to
  node[below, lbl, pos=0.5] {retrieved chunks} (gen.south west);

% Online label
\node[lbl, above=0.15cm of query, font=\tiny\bfseries, gray] {Online Query};

\end{tikzpicture}

  \vspace{0.3em}
  \textbf{Two phases:} (1) Offline indexing — encode documents into a vector store.
  (2) Online query — retrieve, optionally rerank, then generate.
\end{frame}

\begin{frame}{Original RAG Paper}
  \papercite{Retrieval-Augmented Generation for Knowledge-Intensive NLP Tasks}
    {Lewis, Perez, Piktus et al.\ (Facebook AI / UCL)}
    {NeurIPS 2020~\cite{lewis2020retrieval}}
  \vspace{0.4em}
  \textbf{Key contributions:}
  \begin{itemize}
    \item Formalized RAG as a conditional generative model:
          $p(y|x) = \sum_z p_\eta(z|x)\, p_\theta(y|x, z)$
    \item Two variants: RAG-Token (retrieve per token) and RAG-Sequence (retrieve once).
    \item Showed significant gains on open-domain QA (NQ, TriviaQA, WebQ).
    \item Retriever (DPR) and generator (BART) jointly trained.
  \end{itemize}
\end{frame}

\section{Chunking Strategies}
\sectionframe{Chunking}{How you split documents matters enormously}

\begin{frame}{Chunking: The Unsung Hero of RAG}
  \begin{tcolorbox}[alertbox, title=Why Chunking Matters]
    If a chunk is too large → retrieved context diluted; LLM ignores the answer.\\
    If a chunk is too small → answer spans multiple chunks; retrieval misses context.
  \end{tcolorbox}
  \vspace{0.4em}
  \textbf{Chunking strategies:}
  \begin{enumerate}
    \item \textbf{Fixed-size} (e.g., 512 tokens with 50-token overlap): simple, fast.
    \item \textbf{Sentence-based}: split at sentence boundaries.
    \item \textbf{Semantic chunking}: split when embedding similarity drops significantly.
    \item \textbf{Recursive character splitting} (LangChain default): try paragraph → sentence → word.
    \item \textbf{Document-aware}: split by Markdown headers, HTML tags, PDF structure.
    \item \textbf{Small-to-big}: store small chunks, retrieve large parent contexts.
  \end{enumerate}
\end{frame}

\section{Advanced RAG Techniques}
\sectionframe{Advanced RAG}{Query rewriting, reranking, and beyond}

\begin{frame}{Query Rewriting and HyDE}
  \textbf{Problem:} User queries are often terse, ambiguous, or multi-hop.
  \vspace{0.4em}
  \begin{columns}[T]
    \column{0.52\textwidth}
      \textbf{Query rewriting:}
      \begin{itemize}
        \item Rephrase the query before retrieval.
        \item Expand acronyms, disambiguate terms.
        \item Multi-query: generate $k$ query variants; retrieve for each; merge.
        \item Step-back prompting: generalize the query first.
      \end{itemize}
      \vspace{0.4em}
      \textbf{HyDE} (Gao et al.\ 2022):
      \begin{itemize}
        \item LLM generates a hypothetical answer document.
        \item Embed \emph{the hypothetical document} for retrieval.
        \item Bridges query-document vocabulary gap.
      \end{itemize}
    \column{0.44\textwidth}
      \placeholder{0.95\textwidth}{4cm}{HyDE: query → hypothetical doc → embedding → retrieve}
  \end{columns}
\end{frame}

\begin{frame}{Reranking}
  \textbf{Problem:} Bi-encoder retrieval is fast but imprecise.
  \vspace{0.4em}
  \begin{tcolorbox}[hitbox, title=Two-Stage Retrieval]
    \textbf{Stage 1 (retrieve):} Bi-encoder ANN search, return top-100 candidates.\\
    \textbf{Stage 2 (rerank):} Cross-encoder reads query + candidate jointly; reorder top-$k$.
  \end{tcolorbox}
  \vspace{0.4em}
  \textbf{Reranking models:}
  \begin{itemize}
    \item \textbf{Cross-encoder} (e.g., \texttt{ms-marco-MiniLM-L-6-v2}): slow but accurate.
    \item \textbf{Cohere Rerank, Jina Rerank}: API-based, state-of-the-art.
    \item \textbf{LLM-as-reranker}: use LLM to assess relevance (expensive, accurate).
    \item \textbf{RankGPT}: GPT reranking via sliding window.
  \end{itemize}
\end{frame}

\section{Self-RAG}
\sectionframe{Self-RAG}{The model decides when and how to retrieve}

\begin{frame}{Self-RAG: Selective, Self-Reflective RAG}
  \textbf{Asai et al.\ (2023)}~\cite{asai2023selfrag}
  \vspace{0.4em}
  \begin{tcolorbox}[hitbox, title=Self-RAG Tokens]
    Augment the vocabulary with special \emph{reflection tokens}:
    \begin{itemize}
      \item \texttt{[Retrieve]}: should I retrieve now?
      \item \texttt{[Relevant]} / \texttt{[Irrelevant]}: is this passage on-topic?
      \item \texttt{[Supported]} / \texttt{[Partially Supported]} / \texttt{[Not Supported]}: is the generation grounded?
      \item \texttt{[Utility]}: is the response useful?
    \end{itemize}
  \end{tcolorbox}
  \vspace{0.3em}
  The model is trained end-to-end to generate these tokens and act accordingly.
  Retrieval becomes \textbf{conditional}: only triggered when the model predicts \texttt{[Retrieve]=True}.
\end{frame}

\section{Graph RAG}
\sectionframe{Graph RAG}{Querying structured knowledge}

\begin{frame}{Graph RAG: From Local to Global}
  \textbf{Edge et al.\ (2024)}~\cite{edge2024graphrag} — Microsoft Research.
  \vspace{0.4em}
  \begin{columns}[T]
    \column{0.55\textwidth}
      \textbf{Motivation:} Flat chunk retrieval fails for:
      \begin{itemize}
        \item Multi-hop questions (entity A $\to$ entity B $\to$ entity C)
        \item Global summarization (``What are the main themes of this corpus?'')
        \item Cross-document entity resolution
      \end{itemize}
      \vspace{0.4em}
      \textbf{Graph RAG pipeline:}
      \begin{enumerate}
        \item Extract entities + relations from chunks (LLM-based).
        \item Build a knowledge graph.
        \item Detect communities in the graph.
        \item Generate community summaries.
        \item Answer global queries by traversing summaries.
      \end{enumerate}
    \column{0.41\textwidth}
      \placeholder{0.95\textwidth}{4cm}{Knowledge graph + community summaries for global QA}
  \end{columns}
\end{frame}

\section{RAG Evaluation}
\sectionframe{RAG Evaluation}{How do we know if RAG is working?}

\begin{frame}{Evaluation Dimensions}
  \begin{tcolorbox}[hitbox, title=RAGAS Framework~\cite{es2023ragas}]
    \begin{itemize}
      \item \textbf{Faithfulness}: Does the answer follow from the retrieved context?
            Score: entailment between answer sentences and context.
      \item \textbf{Answer Relevance}: Is the answer relevant to the question?
            Score: reverse-generate questions from the answer; measure similarity.
      \item \textbf{Context Precision}: Are retrieved chunks relevant? (precision@k)
      \item \textbf{Context Recall}: Were all necessary facts retrieved? (recall@k)
    \end{itemize}
  \end{tcolorbox}
  \vspace{0.3em}
  RAGAS is \textbf{reference-free}: it uses an LLM judge, not ground-truth answers.
\end{frame}

\begin{frame}{Common RAG Failure Modes}
  \begin{center}
  \begin{tabular}{lll}
    \toprule
    \textbf{Failure} & \textbf{Cause} & \textbf{Fix} \\
    \midrule
    Retrieves wrong chunks & Poor embeddings or chunking & Better splitter, reranker \\
    LLM ignores context & Context too long / irrelevant & Improve retrieval precision \\
    LLM hallucinates anyway & Insufficient context & Add more chunks \\
    Multi-hop fails & Flat retrieval & Graph RAG, iterative retrieval \\
    Outdated answers & Index not refreshed & Incremental indexing \\
    \bottomrule
  \end{tabular}
  \end{center}
\end{frame}

\section{Lab / Demo}
\begin{frame}{Lab Idea: Build and Evaluate a RAG Pipeline}
  \begin{tcolorbox}[hitbox, title=Lab 6 — End-to-End RAG]
    \textbf{Goal:} Build a complete RAG system and evaluate with RAGAS.\\[4pt]
    \textbf{Tools:} LlamaIndex or LangChain, \texttt{sentence-transformers}, Chroma, \texttt{ragas}, GPT-4o-mini
    \begin{enumerate}
      \item Index a set of Wikipedia articles or PDF documents.
      \item Implement a naive RAG (chunk 512 tokens, cosine retrieval, top-3).
      \item Evaluate baseline RAGAS scores (faithfulness, relevance, precision, recall).
      \item Add a reranker (e.g., cross-encoder). Measure RAGAS improvement.
      \item Try HyDE query expansion. Measure RAGAS improvement.
      \item Report: which technique gave the biggest gain?
    \end{enumerate}
  \end{tcolorbox}
\end{frame}

\section{Discussion \& Summary}
\begin{frame}{Discussion Questions}
  \begin{enumerate}
    \item RAG reduces hallucination by grounding answers in retrieved text. But what happens when the retrieved text itself contains false information?
    \item Self-RAG adds reflection tokens to the model vocabulary. What are the training data requirements? Could you apply this to a commercial API model?
    \item Graph RAG builds a knowledge graph before retrieval. When is the cost of knowledge graph construction justified? When is it not?
    \item RAGAS uses an LLM to evaluate an LLM. What are the circular evaluation risks? How would you ground RAG evaluation in human annotation?
  \end{enumerate}
\end{frame}

\begin{frame}{Summary}
  \begin{itemize}
    \item \textbf{RAG} grounds LLM generation in retrieved external documents, reducing hallucination.
    \item \textbf{Chunking} strategy significantly affects retrieval quality; semantic and small-to-big strategies outperform naive fixed-size.
    \item \textbf{Query rewriting} (HyDE, multi-query) bridges vocabulary gaps.
    \item \textbf{Reranking} (cross-encoder) improves precision at modest latency cost.
    \item \textbf{Self-RAG} makes retrieval conditional and self-critiquing.
    \item \textbf{Graph RAG} handles multi-hop and global queries via knowledge graphs.
    \item \textbf{RAGAS} provides reference-free LLM-based evaluation of RAG pipelines.
  \end{itemize}
  \vspace{0.3em}
  \textbf{Next week:} From retrieval to agency — foundations of AI agents.
\end{frame}

\begin{frame}[allowframebreaks]{Suggested Reading}
  \nocite{lewis2020retrieval, asai2023selfrag, edge2024graphrag, es2023ragas,
          karpukhin2020dpr, khattab2020colbert}
  \bibliography{../../references}
\end{frame}

\end{document}
